% Generated draft: requires mathematical and source review.
\documentclass[11pt]{article}
\usepackage[margin=25mm]{geometry}
\usepackage{fontspec}
\setmainfont{Latin Modern Roman}
\usepackage{amsmath,amssymb,mathtools,unicode-math}
\setmathfont{Latin Modern Math}
\usepackage{enumitem,needspace,xurl,hyperref}
\setlength{\parindent}{0pt}
\setlength{\parskip}{5pt}
\setlength{\emergencystretch}{3em}
\newcommand{\N}{\mathbb{N}}
\newcommand{\Z}{\mathbb{Z}}
\newcommand{\Q}{\mathbb{Q}}
\newcommand{\R}{\mathbb{R}}
\newcommand{\C}{\mathbb{C}}
\newcommand{\F}{\mathbb{F}}
\newcommand{\A}{\mathbb{A}}
\newcommand{\PP}{\mathbb{P}}
\newcommand{\E}{\mathbb{E}}
\newcommand{\eps}{\varepsilon}
\newcommand{\dd}{\,\mathrm{d}}
\begin{document}
\section*{UnsolvedMath ID: 20000241}
\textbf{Source problem code:} catalogue entry
\subsection{Definitions and mathematical statement}
The damaged source is reconstructed as a Kronecker-type question proposed by Satya Mandal: are curves in affine space \(\A^n_k\) set-theoretic complete intersections? Algebraically, for a field \(k\), \(n\geq2\), \(R=k[x_1,\ldots,x_n]\), and a radical pure one-dimensional ideal \(I\subset R\), ask whether there exist \(n-1\) polynomials with
\[
I=\sqrt{(f_1,\ldots,f_{n-1})}.
\]
Equivalently, is the arithmetical rank \(\operatorname{ara}(I)\), the least number of equations generating \(I\) up to radical, equal to its height \(n-1\)? This formulation uses scheme-theoretic radicals, not only rational-point zero sets. The original does not fix the field or whether curves are integral, smooth or arbitrary reduced pure curves. Its corrupted extraction and label exact do not independently verify this reconstruction.

The attached aid reports a characteristic-free family. For \(0\leq r<n\), let \(Z_{n,r}\) be the union of all coordinate \(r\)-planes in \(\A^n_k\), and let \(e_j\) be the degree-\(j\) elementary symmetric polynomial in the \(x_i\). It reports
\[
I(Z_{n,r})=\sqrt{J_{n,r}},\qquad
J_{n,r}=(e_{r+1},\ldots,e_n).
\]
For a coordinate component with free variables indexed by \(I\), \(|I|=r\), write \(K=k(x_i:i\in I)\), let \(y_1,\ldots,y_{n-r}\) be the transverse variables, and localize at its generic point. The aid reports the presentation
\[
(R/J_{n,r})_{(x_j:j\notin I)}
\simeq
\frac{K[y_1,\ldots,y_{n-r}]_{(y)}}{(e_1(y),\ldots,e_{n-r}(y))},
\]
the type-\(A\) coinvariant algebra with generic length \((n-r)!\). This multiplicity belongs to the chosen complete-intersection thickening, not the reduced union, whose generic multiplicity is one. In particular, for \(r=1\), the coordinate-axis union is a set-theoretic complete intersection over every field; for \(n\geq3\) it is a reduced curve that is not a local complete intersection at the origin. The latter non-lci assertion does not apply to the two-axis plane hypersurface.

The aid also reports, for every radical pure affine curve ideal as above,
\[
\operatorname{cd}(I,R):=\max\{j:H_I^j(R)\ne0\}=n-1.
\]
It concludes that ordinary local cohomology cannot distinguish its remaining arithmetical-rank alternatives \(n-1\) and \(n\). It says the general characteristic-zero question remains open; no general resolution is asserted from these skeleton examples or the cohomological-dimension calculation.
\subsection{Short English statement}
Ask whether every affine curve can be defined set-theoretically by as many equations as its codimension. The attached aid reports explicit coordinate-skeleton examples and their transverse thickening multiplicities, and says ordinary local cohomology leaves the general arithmetical-rank question undecided.
\subsection{Sources}
\begin{itemize}
\item Local UnsolvedMath registry, permanent ID \(20000241\), exact public code \texttt{catalogue entry}; canonical location \texttt{aim-algebraic-geometry-notes.json}, notes[240].
\item AIM workshop ``Projective modules and \(A^1\)-homotopy theory,'' extracted Question 4, with the damaged attribution ``Krone ker-prop osed by Sat ya Mandal''; supplied URL \url{https://aimath.org/pastworkshops/projectiveA1problems.pdf}. The PDF was not consulted in this drafting step. No year or full literature citations are supplied; the literature field is empty.
\item Attached research attempt 1, ``Coordinate-skeleton curves, transverse coinvariant algebras, and the cohomological-dimension gap,'' without complete bibliographic references.
\end{itemize}
\end{document}
